\centerline{\bf \Huge Вспоминаем комбу}

\begin{flushright}
        \scriptsize{\it Я не боюсь узнать свои слабости, поскольку они — часть меня самого.\\ А я принимаю себя таким, какой я есть!\\ Кадзи}
\end{flushright}

\problem{1}
В ЭлМШ 99 учеников. Каждый ботает 7 дней подряд в году, в остальные дни --- отдыхает. Докажите, что число дней, когда ботает нечетное число учеников, не меньше 7.

\problem{2}
Рассмотрим операцию

$$
x \star y=\frac{x+y}{x y+4} .
$$

Найдите значение выражения

$$
0 \star(1 \star(\ldots(2024 \star 2025) \ldots))
$$

\problem{3}
Было 12 карточек с надписями "Слева от меня ровно 1 ложное утверждение"\ , "Слева от меня ровно 2 ложных утверждения"\ , $\ldots$, "Слева от меня ровно 12 ложных утверждений"\ . Даяна разложила карточки в ряд слева направо в каком-то порядке. Какое наибольшее число утверждений могло оказаться истинными?

\problem{4}
В 9в классе 25 учеников. Известно, что у любых двух девочек класса количество друзей-мальчиков из этого класса не совпадает. Какое наибольшее количество девочек может быть в этом классе?

\problem{5}
В какое наибольшее число цветов можно раскрасить доску $8 \times 8$ так, чтобы каждая клетка граничила по стороне хотя бы с двумя клетками своего цвета?

\problem{6}
В корпорации ЭлМШ работают 10002 сотрудника. Все работни пришли на день рождения Расула Владимировича и он их рассадил за один круглый стол. Известно, что зарплаты сидящих рядом различаются на 2 или 3 ЭлМШ коина. Насколько сильно может отличаться зарплата двух работников, если все их зарплаты различны?
%15002
%https://problems.ru/view_problem_details_new.php?id=98590

\problem{7}
На шахматной доске стоят 10 шахматных фигур - слоны и ладьи. Никакая фигура не бьёт ни одну другую. Какое наименьшее количество слонов может быть среди них?


\newpage

\hspace{3mm}

\centerline{\bf \Huge Забываем комбу...}

\begin{flushright}
        \scriptsize{Кто пришёл без сменки, тот лох!!!}
\end{flushright}

\problem{1}
В классе учится 23 человека. В течение года каждый ученик этого класса один раз праздновал день рождения, на который пришли некоторые (хотя бы один, но не все) его одноклассники. Могло ли оказаться, что каждые два ученика этого класса встретились на таких празднованиях одинаковое число раз? (Считается, что на каждом празднике встретились любые два гостя, а также именинник встретился со всеми гостями.)
%[Регион, 2016, 9.5 (из 8)]


\problem{2} 
На окружности отмечены 77 точек, среди них нет диаметрально противоположных. Докажите, что можно провести диаметр через одну из точек так, что по обе стороны диаметра точек окажется поровну.

\problem{3}
В каждой клетке квадрата $2012 \times 2012$ записано вещественное число от 0 до 1. Рассмотрим всевозможные разрезы этого квадрата на два прямоугольника по вертикальной или горизонтальной линии сетки. Оказалось, что какой бы разрез не был произведен, сумма чисел хотя бы в одном из двух получившихся прямоугольников не больше 1. Найдите максимально возможную сумму чисел в исходном квадрате.
%APMO 2012

\problem{4}
Образцово-показательная семья носит одежду трех цветов: белого, синего и красного. Каждую неделю они занашивают 10 килограммов одежды и распределяют ее по трем одинаковым бакам: с белым бельем, с синим бельем и с красным бельем (отношение белого к синему к красному может менять от недели к неделе). В конце недели стирается вся одежда из бака, в котором собралась максимальная масса белья. Какова минимальная вместимость бака, при которой ни при каких обстоятельствах не произойдет переполнения?

\problem{5} 
Петя и Вася играют в следующую игру. Изначально на доске написано число $2011^{2011}$. Они по очереди, начинает Петя, либо вычитают из него число от 1 до 2010 включительно, либо делят на 2011, округляя вниз. Выигрывает тот, кто первым получит неположительное число. Кто выигрывает при правильной игре?

\problem{6} 
В ориентированном графе нет ориентированных циклов, а также путей длины более 99. Докажите, что ребра графа можно раскрасить в два цвета так, чтобы не было одноцветных путей длины 10.